\section{Results}
\label{sec:wins}

DG's latency advantage is dramatic in serial
mode. AR-30B model takes roughly six seconds to produce a
123-token JSON record; DG completes the same task in under 400\,ms.
To make this concrete: if both models serve a security operations
centre processing alerts one at a time, DG handles \crossoverLambda{}
alerts per second per GPU; the AR-30B handles fewer than one in five
seconds. Or, in human terms: DG finishes before you can blink; AR-30B
finishes before you finish reading this sentence. AR stops
at its natural end-of-sequence token (EOS), so latency plateaus regardless of
\texttt{max\_tokens}.

Table~\ref{tab:niche} shows the serial comparison, AR-30B measured,
DG from the staircase law of Eq.~\ref{eq:dg}. The
pattern reveals two things: first, AR latency is flat from 128 tokens
onward because the model hits EOS, setting \texttt{max\_tokens}
higher does nothing. Second, DG latency stays flat for the first
canvas and then steps up per additional canvas. The advantage is
therefore non-monotonic: it climbs to $15.6\times$ once AR reaches its
natural length, then narrows to $8.0\times$ by $L{=}1024$ as DG starts
paying for canvases while AR sits at its plateau. The $4.2\times$ in
the first row is not a like-for-like comparison, at
\texttt{max\_tokens}${=}32$ AR is truncated mid-record while DG still
emits a full canvas. Extending the same accounting past the table, the
ratio keeps falling ($6.9\times$ at $L{=}1280$, $4.9\times$ at
$L{=}2048$), so the eightfold floor holds across the budgets tabulated
here, not universally.

\begin{table}[h]
\centering
\caption{Serial latency at the MoE-matched pair (AR-30B vs DG,
batch-1, BGL). AR-30B column is measured ($n{=}50$) and plateaus
because the model stops at $\sim$123 tokens (EOS), so raising
\texttt{max\_tokens} buys nothing. The DG column is Eq.~\ref{eq:dg}
evaluated at the canvas count each budget implies, anchored on the
measured base ($n{=}300$) and canvas step ($n{=}200$); forcing DG to
fill those canvases measures 501\,ms at $L{=}512$ and 739\,ms at
$L{=}1024$, within 1\% of the law. The two columns therefore compare
AR at its natural output against DG at a full canvas budget, a
conservative reading for DG, which emits four times more tokens at the
last row.}
\label{tab:niche}
\begin{tabular}{@{}lrrrr@{}}
\toprule
max\_tokens & AR-30B (ms) & AR actual out & DG (ms) & DG advantage \\
\midrule
32   & 1604 & 32  & 379 & $4.2\times$ \\
128  & 5923 & 121 & 379 & $15.6\times$ \\
256  & 5899 & 123 & 379 & $15.6\times$ \\
512  & 5898 & 123 & 498 & $11.8\times$ \\
1024 & 5910 & 123 & 737 & $8.0\times$ \\
\bottomrule
\end{tabular}
\end{table}

\paragraph{DG wins at every practical output length.}

The serial cost law gives crossover $L^{*} = \crossoverL$ tokens, DG
is faster at any output length above $L^{*}$ (integer-safe: $L \geq 7$).
AR-30B serial slope (\arSlope\,ms/token) is $73\times$ steeper than
its batch-throughput slope (0.65\,ms/token at batch 300, from
80.0\,ms/event over 123 tokens), because serial mode pays full prefill
and decode overhead per request while batch mode amortises across
concurrent requests. Quoting a batch-throughput slope as if it were a
latency therefore understates AR serial latency by a factor of
seventy.

\paragraph{Energy.}

DG uses less energy per event because it draws far
less power above idle: 10.8\,W net against the AR-30B's 69.1\,W. It is
not that DG is busy for less time, in these two runs it was busy
longer per event, but that denoising a canvas keeps fewer units fed
than driving a token-at-a-time decode loop. Over a million events the
gap is 1.6\,kWh, about an hour of a space heater. The energy ratio is:
\begin{equation}
R_{\mathrm{energy}} = \frac{E_{\mathrm{AR}}}{E_{\mathrm{DG}}}
\end{equation}
where $E_{\mathrm{AR}}$ and $E_{\mathrm{DG}}$ are the net energy per
event (joules, idle-subtracted) for each arm. Measured values:
\begin{align}
\text{AR-30B MoE:} \quad & R_{\mathrm{energy}} =
\frac{10.1}{4.2} = 2.4 \quad \text{(idle 92.7\,W)} \\
\text{Dense AR-1.5B:} \quad & R_{\mathrm{energy}} =
\frac{5.1}{4.2} = 1.2 \quad \text{(idle 58.6\,W)}
\end{align}
A ratio above 1 means AR uses more energy. Two caveats make this
indicative instead of controlled: the three runs used different idle
baselines and the DG energy was never re-measured at the AR-30B idle
point; and the AR-30B energy run was served in batch (147\,ms/event
wall time) while the DG run was serial (386\,ms/event), so the two
arms were not in the same serving regime.

\paragraph{Graceful degradation under deadline pressure.}

Under a tightened deadline the diffusion model returns a complete canvas
at lower quality, while the AR model returns a truncated record that
fails schema validation. The $K$ knob is the difference: drop $K$ from
8 to 4 and the model still fills every field, just with slightly
less precision, instead of cutting the JSON in half mid-key. At a
p99 budget of 400\,ms, this meets it at \kFourPninetynine\,ms p99 with
F1 falling from \kEightFone\ to \kFourFone. Table~\ref{tab:contrast}
summarises the contrast across practical challenges.

\begin{table}[h]
\centering
\caption{Practical challenge contrast. Each row is a deployment
scenario; the right column names the architectural property that drives
the outcome. Every AR cell is the AR-30B MoE except the tail-latency
statistics, which are the dense AR-3B, the only arm measured
per-event. Because that arm is dense, the tail asymmetry can't be
attributed to MoE routing; it's DG's denoising-step variance that
drives it, and the AR-30B MoE's own tail is unmeasured.}
\label{tab:contrast}
\begin{tabularx}{\linewidth}{@{}L L L L@{}}
\toprule
\thead{Challenge} & \thead{AR-30B MoE} & \thead{DG-26B MoE} & \thead{Key driver} \\
\midrule
Long output, batch-1 & \arSlope\,ms/token (serial) & Staircase, \dgStep\,ms/canvas & Crossover at $L^{*}\!\approx\!\crossoverL$ ($L\geq7$) \\
Deadline pressure & Truncates output & Drops $K$; complete canvas & Per-request quality knob \\
High arrival rate & Batches $77\times$ & Caps at $\leq$1.7$\times$ & KV-cache reuse \\
Tail-sensitive service-level agreement (SLA) & Coefficient of variation (CV) \arEventCV\%, max/med \arMaxMed$\times$ (dense AR-3B) & CV \dgEventCV\%, max/med \dgMaxMed$\times$ & Denoising-step variance \\
Energy per event & 10.1\,J (batch-served) & 4.2\,J (serial) & Net power draw \\
Long prompt & Prefill paid once & Prefill paid $K$ times & Re-attention per step \\
\bottomrule
\end{tabularx}
\end{table}

